\definecolor{BaselineBlue}{HTML}{E3F5FB}
\definecolor{EvoGainGreen}{RGB}{0,140,45}
\definecolor{EvoLossRed}{RGB}{190,20,32}
\definecolor{EvoNeutralGray}{RGB}{110,110,110}
\newlength{\EvoDeltaWidth}
\settowidth{\EvoDeltaWidth}{{\scriptsize +00.0}}


\newcommand{\gain}[1]{\,\textcolor{EvoGainGreen}{\scriptsize +#1}}
\newcommand{\loss}[1]{\,\textcolor{EvoLossRed}{\scriptsize -#1}}
\newcommand{\same}{\,\textcolor{EvoNeutralGray}{\scriptsize $\pm$0.0}}
\newcommand{\nodelta}{%
  \,\makebox[\EvoDeltaWidth][l]{}%
}

\begin{table*}[t]
\centering
\caption{
Evo-Bench harness-evolving capability leaderboard.
Methods are ranked by Overall score; the two baseline harnesses are excluded from ranking. \textbf{Bold} and \underline{underlined} denote the best and second-best raw scores for each metric, respectively. 
\textcolor{EvoGainGreen}{Green} and \textcolor{EvoLossRed}{red} values indicate absolute score changes relative to the CodeAct baseline; \textcolor{EvoNeutralGray}{gray} indicates no change.
}
\label{tab:main-leaderboard}
% \setlength{\tabcolsep}{5pt}
\small
\resizebox{\textwidth}{!}{%
\begin{tabular}{c l ccccc}
\toprule
\textbf{\rule{0pt}{2ex}Rank}
& \textbf{Model}
& \textbf{Search score}
& \textbf{Office score}
& \textbf{General score}
& \textbf{Overall score}
& \textbf{AnytimeVal} \\
% [2pt]
\midrule
\rowcolor{BaselineBlue}
\multicolumn{7}{c}{\textit{Frontier models}} \\
\midrule
1 & GPT-5.6 Sol~\cite{gpt56}
  & 44.5\gain{32.8}
  & \underline{41.6}\gain{3.2~}
  & \textbf{59.4}\gain{11.0}
  & \textbf{46.3}\gain{16.6}
  & 50.1 \\
2 & Claude Opus 4.8~\cite{anthropic2026opus}
  & \textbf{46.5}\gain{34.8}
  & 39.7\gain{1.3~}
  & \underline{56.3}\gain{7.9~}
  & \underline{45.8}\gain{16.1}
  & \textbf{51.4} \\
3 & GLM-5.2~\cite{glm52}
  & \underline{45.4}\gain{33.7}
  & 39.2\gain{0.8~}
  & 48.4\same
  & 43.5\gain{13.8}
  & \underline{51.0} \\
4 & Qwen3.7-Max~\cite{qwen37}
  & 36.3\gain{24.6}
  & 37.8\loss{0.6~}
  & \textbf{59.4}\gain{11.0}
  & 41.5\gain{11.8}
  & 49.3 \\
5 & Minimax-M3~\cite{lai2026minimax}
  & 33.6\gain{21.9}
  & \textbf{41.7}\gain{3.3~}
  & \underline{56.3}\gain{7.9~}
  & 41.4\gain{11.7}
  & 49.0 \\
7 & Deepseek V4 Pro~\cite{deepseekv4}
  & 34.4\gain{22.7}
  & 39.1\gain{0.7~}
  & 48.4\same
  & 39.1\gain{9.4~}
  & 45.4 \\
8 & Kimi K2.7 Code~\cite{kimik27}
  & 34.5\gain{22.8}
  & 38.1\loss{0.3~}
  & 48.4\same
  & 38.7\gain{9.0~}
  & 43.4 \\
\midrule
\rowcolor{BaselineBlue}
\multicolumn{7}{c}{\textit{Open-weight models}} \\
\midrule
6 & Qwen3.6-27b~\cite{qwen3.6-27b}
  & 34.8\gain{23.1}
  & 38.8\gain{0.4}
  & 50.0\gain{1.6}
  & 39.4\gain{9.7}
  & 46.9 \\
9 & Gemma-4-31B~\cite{gemmateam2026gemma4}
  & 24.2\gain{12.5}
  & 40.4\gain{2.0}
  & 50.0\gain{1.6}
  & 35.9\gain{6.2}
  & 36.2 \\
\midrule
\rowcolor{BaselineBlue}
\multicolumn{7}{c}{\textit{Baseline variants}} \\
\midrule
-- & CodeAct
   & 11.7\nodelta
   & ~38.4\nodelta
   & ~48.4\nodelta
   & ~29.7\nodelta
   & -- \\
-- & Artificial Harness
   & 46.7\nodelta
   & ~43.9\nodelta
   & ~56.3\nodelta
   & ~47.5\nodelta
   & -- \\
\bottomrule
\end{tabular}
}
\end{table*}



\section{Experiments}
\label{sec:experiments}
\raggedbottom

% In this section, we present the main results and discuss in which cases agents outperform official instruction-tuned models.

% 在这个section， we present the main results and并且讨论在什么情况下，模型进化出的框架能够超过现在的最佳框架。
In this section, we present the main results of Evo-Bench. We evaluate the harness-evolving capabilities of frontier models, analyze their evolve process, and further discuss under what conditions model-evolved harnesses can surpass the current state-of-the-art human-engineered frameworks.

\subsection{Experimental Setup}
\label{sec:exp-setup}

\paratitle{Baselines.}
We compare the evolved harnesses against two baselines: (1) the initial \textbf{CodeAct} seed harness, which serves as the starting point to quantify absolute evolutionary gains, and (2) \textbf{Artificial harness}, a composite of domain-specific human-engineered frameworks: MiroFlow~\cite{miromind2025mirothinker} for search, Stirrup~\cite{artificialanalysis2026stirrup} for office, and Claw-Eval~\cite{claweval} for general tasks.
% Other details follow:


\subsubsection{Evaluation protocol.}
Unless stated otherwise, all experiments use DeepSeek-V4-Flash as the fixed policy model and start from the same CodeAct seed policy harness $H_0$. The main experiments use a common budget of 20 iterations, 1,000 evolver steps, and 48 hours. Each policy rollout is capped at 300 steps and one hour. Search and office tasks use one rollout, whereas Claw-Eval use three rollouts, following its native three-trial $\mathrm{Pass}$\^{}$3$ metric. Tasks requiring LLM-based grading use Qwen3.7-Plus as the common judge. We run all experiments one time, and report domain-level scores, Overall Score, and Anytime Validation Score (AnytimeVal), as defined in the previous section. Costs are computed from recorded input, cached-input, and output tokens using provider prices as of July 10, 2026.

\subsubsection{Model configuration.}
We evaluate seven frontier models and two open-weight models on Evo-Bench across the following settings. 
All evolvers run in thinking mode with temperature $1.0$ and the largest context window supported by the model. For models exposing a reasoning-effort control, we use their maximum setting; for models that do not expose a configurable reasoning-effort level, we use their default setting. The shared DeepSeek-V4-Flash policy model uses maximum reasoning effort, temperature $1.0$, and a 256K-token context window. The Qwen3.7-Plus judge uses temperature $0.0$ and a 1M-token context window. 
% We evaluate seven frontier evolvers---Qwen3.7-Max, MiniMax M3, DeepSeek-V4-Pro, Kimi-K2.7-Code, GLM-5.2, GPT-5.6-Sol, and Claude Opus 4.8---and two open-weight evolvers, Qwen3.6-27B and Gemma 4 31B. 
% : \texttt{max} for GPT, Claude, GLM, and DeepSeek. Qwen3.7-Max, MiniMax M3, Qwen3.6-27B, Gemma 4 31B, and Kimi-K2.7-Code
% Exact model identifiers, providers, and inference configurations are listed in Appendix~\ref{app}.

\subsection{Main Results}


\paratitle{Overall performance.}
Table~\ref{tab:main-leaderboard} presents the overall leaderboard of all evaluated models. GPT-5.6-Sol and Claude Opus-4.8 lead with scores of 46.3 and 45.8 respectively, yielding massive absolute gains over the initial CodeAct baseline. This widespread positive delta validates that frontier LLMs possess a genuine capability to autonomously optimize executable harnesses. Examining evolutionary progress, Claude Opus-4.8 and GLM-5.2 achieve the highest Anytime Validation scores. This early peak suggests they rapidly evolve high-quality structures but frequently introduce detrimental modifications in subsequent iterations. Despite these gains, the top evolved harness still slightly underperforms domain-specific Artificial harness composite score of 47.5, leaving headroom for future research.

% Table~\ref{tab:main-leaderboard} presents the overall performance leaderboard on Evo-Bench. GPT-5.6-Sol and Claude Opus-4.8 demonstrate leading capabilities, securing top Overall scores of 46.3 and 45.8 respectively. Notably, all evaluated models achieve substantial improvements over the initial CodeAct baseline score of 29.7. For instance, GPT-5.6-Sol yields a massive absolute overall gain of +16.6, while the open-weight Qwen3.6-27B achieves a solid +9.7 gain, outperforming several proprietary frontier models. This widespread positive delta clearly validates that current large language models possess a genuine, intrinsic capability to autonomously optimize executable agent harnesses.

% Examining the temporal dynamics of this evolutionary process reveals further insights into model behavior. Claude Opus-4.8 and GLM-5.2 achieve the highest Anytime Validation scores of 51.4 and 51.0 respectively. This early peak suggests that while these models rapidly evolve high-quality structures in initial stages, they frequently introduce detrimental modifications during subsequent iterations that degrade performance. Furthermore, despite these remarkable autonomous gains, a performance gap remains when comparing the final results to carefully designed artificial harness. The top-performing evolved harness from GPT-5.6-Sol still slightly underperforms the domain-specific Official harness composite score of 47.5. This gap highlights that fully autonomous harness engineering has not yet completely reached the pinnacle of manual optimization, leaving vital headroom for future agentic self-improvement research.

\paratitle{Per-domain performance.}
The evolutionary gains are highly uneven across different task domains. \textbf{Search} tasks are more amenable to optimization, with Claude Opus-4.8 gaining +34.8 to nearly match the Artificial harness, indicating evolvers easily synthesize missing web-navigation logic. Conversely, \textbf{Office} tasks remain stubborn. Most models show marginal improvements or slight regressions, failing to match the Artificial baseline likely due to the difficulty of discovering specialized workflows. Interestingly, for \textbf{General} tasks, top evolvers like GPT-5.6-Sol and Qwen3.7-Max strictly surpass the Artificial harness. This milestone demonstrates that autonomously evolved reasoning structures can indeed outstrip manually crafted solutions.

% The evolutionary gains are highly uneven across different task domains. Models exhibit the most dramatic improvements in the \textbf{Search} domain. For example, Claude Opus-4.8 achieves a remarkable absolute gain of +34.8 over the baseline, reaching a score of 46.5 that nearly closes the gap with the Official harness score of 46.7. This massive leap suggests the initial CodeAct loop is fundamentally deficient in search-specific capabilities like HTML parsing, which frontier evolvers can successfully synthesize from scratch. Conversely, performance on \textbf{Office} tasks remains stubborn and still falls noticeably short of the Official harness score of 43.9. Models show only marginal improvements in this area, with a peak gain of +3.3 for Minimax M3, while others like Qwen3.7-Max experience a slight regression of -0.6. This difficulty implies that Office environments likely require rigid and highly specialized API integrations that are challenging to discover through general evolution. Interestingly, for \textbf{General} agent tasks, top evolvers demonstrate exceptional capability. Both GPT-5.6-Sol and Qwen3.7-Max achieve a score of 59.4, strictly surpassing the Official harness score of 56.3. This milestone demonstrates that autonomously evolved reasoning structures can sometimes outstrip manually crafted solutions.

\paratitle{Budget use.}
Figure~\ref{fig:budget_use} reveals two distinct evolutionary behaviors: exhaustive exploration and early saturation. GPT-5.6-Sol and Kimi-K2.7-Code uniquely exhaust the maximum 20-iteration budget with the highest step counts and longest durations. This persistent exploration directly contributes to the top-tier overall performance of GPT-5.6-Sol. Conversely, most models terminate prematurely due to invalid code proposals or stagnant reasoning loops. For instance, Qwen3.7-Max and DeepSeek-V4-Pro halt at 15 iterations using barely 200 steps. However, budget exhaustion does not strictly dictate success. Despite its early halt, Qwen3.7-Max delivers a highly competitive score and dominates the General tasks domain, demonstrating exceptional sample efficiency in synthesizing effective updates without exhaustive trial and error.

\begin{figure}[htbp]
    \centering
    \includegraphics[width=0.9\linewidth]{Figures/budget_use.pdf}
    \caption{Budget usage of each evolver. Iters, steps and time report the number used relative to the maximum budgets of 20 iterations, 1,000 steps, and 48 hours, respectively.}
    \label{fig:budget_use}
\end{figure}

% \begin{figure}[htbp]
%     \centering

%     \begin{subfigure}[t]{0.48\linewidth}
%         \centering
%         \includegraphics[width=\linewidth]{Figures/budget_use.pdf}
%         \caption{Budget usage of each evolver. Iters, steps and time report the number used relative to the maximum budgets of 20 iterations, 1,000 steps, and 48 hours, respectively.}
%         \label{fig:budget_use}
%     \end{subfigure}
%     \hfill
%     \begin{subfigure}[t]{0.48\linewidth}
%         \centering
%         \includegraphics[width=\linewidth]{Figures/evolve_cost_vs_performance.pdf}
%         \caption{Cost(\$) vs. Performance. The dashed line denotes the Pareto frontier of cost-effectiveness, while the horizontal line indicates the performance of the artificial harness.}
%         \label{fig:evolve-cost}
%     \end{subfigure}
%     \caption{}
%     \label{fig:two-images}
% \end{figure}

\paratitle{Cost analysis.}
Figure~\ref{fig:evolve-cost} illustrates the overall performance against the API inference cost incurred exclusively by the evolver model. The Pareto frontier reveals a steep logarithmic trade-off between financial investment and evolutionary capability. Achieving top-tier results demands substantial resources: GPT-5.6-Sol secures the highest score but dominates the cost spectrum by exceeding 500 USD per run. In stark contrast, models occupying the middle ground demonstrate exceptional cost-effectiveness. GLM-5.2 and Qwen3.7-Max establish the highly efficient knee of the curve, delivering formidable performance for under 40 USD. At the extreme budget end, DeepSeek-V4-Pro anchors the absolute baseline, successfully synthesizing functional harness improvements for less than a single dollar.

% Figure~\ref{fig:evolve-cost} illustrates the overall performance against the API inference cost incurred exclusively by the evolver model, intentionally omitting the evaluation expenses. The resulting Pareto frontier reveals a steep logarithmic trade-off between financial investment and evolutionary capability. Achieving top-tier results demands substantial resources. For instance, GPT-5.6-Sol secures the highest score but dominates the cost spectrum by exceeding 500 USD per run, while Claude Opus-4.8 requires a similar massive budget but falls slightly below the optimal frontier. In stark contrast, models occupying the middle ground demonstrate exceptional cost-effectiveness. GLM-5.2 and Qwen3.7-Max establish the highly efficient knee of the curve, delivering formidable performance for under 40 USD. At the extreme budget end, DeepSeek-V4-Pro anchors the absolute baseline of the Pareto frontier, successfully synthesizing functional harness improvements for less than a single dollar. 

\subsection{Further Analysis: When Can Evolved Harnesses Surpass Human-Engineered Systems?}
\label{sec:analysis}

\paratitle{Case study: how GPT-5.6-Sol evolves a harness.}
% Figure~\ref{fig:gpt56-case-study} illustrates GPT-5.6-Sol acting as an autonomous systems engineer by architecting a hierarchical, domain-specific router for prompts and tools rather than a uniform harness. Beyond structural design, it demonstrates sophisticated tool engineering and self-correction. First, for Search, it builds web navigation tools from scratch and adds output sanitation to resolve execution violations. Second, for Office, it engineers specialized workflows like source ledgers while safely reverting detrimental gating mechanisms. Third, for General tasks, it uses regression diffs to diagnose prompt over-steering, restoring a minimal structure while retaining essential safeguards. This refinement explains our findings: while removing interference allows the evolved harness to specifically outperform manual solutions on General tasks, the overall system still trails the public artificial harness. Human experts maintain a crucial edge in Search and Office by providing highly customized infrastructure.
Figure~\ref{fig:gpt56-case-study} shows GPT-5.6-Sol acting as an autonomous systems engineer: instead of applying a uniform harness, it constructs a hierarchical router for domain-specific prompts and tools. It further demonstrates tool construction, failure diagnosis, and iterative correction. For Search, it implements web-search/fetch tools and a webpage cleaner that removes scripts, styles, and markup noise while preserving links. For Office, it separates APEX and GDPval: APEX tracks evidence by file, page, and sheet cell and recomputes key values, whereas GDPval requires artifact creation, recalculation, reopening, and rendered inspection. For General tasks, it adds recovery from empty or premature responses, credential redaction, and a guard against sending draft-only emails. Across evolvers, Opus-4.8 exhibits the closest systems-engineering profile to GPT-5.6-Sol, combining tool construction with domain-aware control, context management, failure recovery, and evidence-based reversion. GLM-5.2 demonstrates substantial diagnosis and tool engineering but relies more on domain-specific guidance and local controls, while Qwen3.7-Max and MiniMax M3 focus primarily on web tools, prompt specialization, and safety filters.


\begin{figure}[htbp]
    \centering
    \begin{minipage}[t]{0.54\linewidth}
        \vspace{0pt}
        \centering
        \includegraphics[width=\linewidth]{Figures/evolve_cost_vs_performance.pdf}
        \caption{Cost(\$) vs. performance. The dashed line denotes the Pareto
        frontier; the horizontal line denotes the artificial harness.}
        \label{fig:evolve-cost}
    \end{minipage}
    \hfill
    \begin{minipage}[t]{0.43\linewidth}
        \vspace{0pt}
        \centering
        \includegraphics[width=\linewidth]{Figures/gpt56_case_study.pdf}
        \caption{GPT-5.6-Sol harness-evolution case study.}
        \label{fig:gpt56-case-study}
    \end{minipage}
\end{figure}
% learns a two-level task router, adds capability for Search, specializes task contracts for Office, and removes prompt interference for General/Claw. Intermediate scores are from validation; the final block reports held-out evaluation.

\paratitle{Challenges.}
% 这里写挑战
Despite its achievements, GPT-5.6-Sol exposes three critical limitations. First, it reacts superficially to aggregate scores rather than distilling causal failure modes from extensive logs. Second, it resolves cross-domain interference via naive domain routing instead of discovering fundamentally robust, shared mechanisms. Third, it underutilizes the research budget, lacking the goal-oriented persistence of human engineers. Consequently, the evolved harness relies on localized modifications, leaving core components primitive: the planner acts passively, the context remains append-only, and the verifier stays entirely permissive. These behavioral patterns extend across other evolvers. Opus-4.8 exhibits the closest systems-engineering profile to GPT-5.6-Sol. In contrast, GLM-5.2 relies more heavily on local controls, while Qwen3.7-Max and MiniMax M3 focus primarily on basic web tools and prompt specialization.

\paratitle{Discussion.}
In summary, current frontier models demonstrate highly promising autonomous engineering capabilities, notably outperforming manual designs on General tasks. However, the overall lack of deep system-level optimization explains why they still trail the Artificial harness on Search with a score of 44.5 against 46.7, and on Office with 41.6 against 43.9. We believe that by strengthening the capacity to abstract causal failure patterns, execute large-scale architectural refactoring, and maximize budgets utilization, future language models hold substantial potential to continuously evolve harnesses that consistently surpass human experts.



% \paratitle{Hacking behaviors.}
% % 这里写一些hack的行为。
